\chapter{Busoni, Ferruccio} % (fold)
\label{cha:busoni}

\href{https://en.wikipedia.org/wiki/Ferruccio_Busoni}{Biography on Wikipedia}

\href{https://en.wikipedia.org/wiki/List_of_compositions_by_Ferruccio_Busoni}{List of Works}

\section{Orchestral}

\noindent{Geharnischte Suite Op. 34a, for orchestra (Second Orchestral Suite) (1894-1895, rev 1903), BV 242} \textit{BBC PO, Neeme J\"{a}rvi}\marginnote{max}

\noindent{Violin Concerto in D major, Op. 35a (1896-1897), BV 243} \textit{Francesca Dego, BBC SO, Dalia Stasevska}\marginnote{max}

\noindent{Eine Lustspielouvertüre (Comedy Overture) Op. 38, for orchestra (1897, rev 1904), BV 245} \textit{BBC PO, Neeme J\"{a}rvi}\marginnote{high}

\noindent{Piano Concerto in C major, Op. 39 (1901-1904), BV 247} \textit{John Ogdon, Royal PO, Daniell Revenaugh}\marginnote{high}

\noindent{Turandot Suite Op. 41, for orchestra (1905), BV 249} \textit{Cincinnati SO, Michael Gielen}\marginnote{high}

\noindent{Berceuse \'{e}l\'{e}giaque, Op.42, for orchestra (1909), BV 252a} \textit{BBC PO, Neeme J\"{a}rvi}\marginnote{max}

\noindent{Die Brautwahl Suite Op. 45, for orchestra (1912), BV 261} \textit{BBC PO, Neeme J\"{a}rvi}\marginnote{high}

\noindent{Nocturne symphonique Op. 43, for orchestra (1912-1913), BV 262} \textit{Staatskapelle Dresden, Christian Thielemann}\marginnote{high}

\noindent{Indianische Fantasie [Indian Fantasy] Op. 44, for piano and orchestra (1913-1914), BV 264} \textit{Nelson Goerner, BBC PO, Neeme J\"{a}rvi}\marginnote{high}

\noindent{Rondò arlecchinesco Op. 46, for orchestra and tenor solo (1915), BV 266} \textit{NDR SO, Werner Andreas Albert}\marginnote{high}

\noindent{Gesang vom Reigen der Geister [Song of Dancing of the Spirits]. Indianisches Tagebuch. Zweites Buch. [Indian Diary. Second Book] Op. 47, for small orchestra (1915), BV 269} \textit{BBC PO, Neeme J\"{a}rvi}\marginnote{high}

\noindent{Concertino Op. 48, for clarinet and small orchestra (1918), BV 276} \textit{John Bradbury, BBC PO, Neeme J\"{a}rvi}\marginnote{max}

\noindent{Sarabande und Cortège Op. 51. Two studies for Doktor Faust, [Fifth and Sixth Elegies for Orchestra] (1818-1819), BV 282} \textit{BBC PO, Neeme J\"{a}rvi}\marginnote{max}

\noindent{Tanzwalzer Op. 53, for orchestra (1920), BV 288} \textit{BBC PO, Neeme J\"{a}rvi}\marginnote{max}

\section{Chamber}

\noindent{Suite Op. 10, for clarinet and piano (1878), BV 88} \textit{Dieter Kl\"{o}cker, Consortium Classicum}\marginnote{high}

\noindent{Solo dramatique (in B-flat minor) Op. 13 (Op. 33), for clarinet and piano (1879), BV 101} \textit{Dieter Kl\"{o}cker, Consortium Classicum}\marginnote{high}

\noindent{Sonata in D major, for piano and clarinet, BV 138} \textit{Dieter Kl\"{o}cker, Consortium Classicum}\marginnote{high}

\noindent{Duo (in E minor) Op. 43 (Op. 73), for 2 flutes and piano (1880), BV 156} \textit{Consortium Classicum}\marginnote{high}

\noindent{Suite (in G minor), for clarinet and string quartet (1880?), BV 176} \textit{Dieter Kl\"{o}cker, Consortium Classicum}\marginnote{high}

\noindent{Kleine Suite Op. 23, for cello and piano (1886), BV 215} \textit{Ludovica Rana, Maddalena Giacopuzzi}\marginnote{high}

\noindent{Sonata (no. 1 in E minor) Op. 29, for violin and piano (1889), BV 234} \textit{Ingolf Turban, Ilja Scheps}\marginnote{high}

\noindent{Kultaselle. Zehn kurze Variationen über ein finnisches Volkslied [Ten short variations on a Finnish Folk Song], for piano and cello (1889), BV 237} \textit{Chiara Burattini, Umberto Jacopo Laureti}\marginnote{max}

\noindent{Sonata no.2 (in E minor) Op. 36a, for violin and piano (1898, 1900), BV 244} \textit{Ingolf Turban, Ilja Scheps}\marginnote{high}

\noindent{Albumblatt [Albumleaf] (in E minor), for flute (or muted violin) and piano (1916), BV 272} \textit{Chiara Burattini, Umberto Jacopo Laureti}\marginnote{max}

\noindent{Elegie (in E-flat major), for clarinet and piano (1919-1920), BV 286}  \textit{Dieter Kl\"{o}cker, Consortium Classicum}\marginnote{high}

\section{Piano}

\noindent{Elegien. 6 neue Klavierstücke. (1907), BV 249} \textit{Peter Donohoe}\marginnote{max}

\noindent{Nuit de Noël. Esquisse [Christmas. Sketch], for piano (1908), BV 251} \textit{Victor Nicoara}\marginnote{high}

\noindent{Berceuse (Lullaby), for piano (1909), BV 252} \textit{Peter Donohoe}\marginnote{max}

\noindent{Fantasia nach Johann Sebastian Bach, for piano (1909), BV 253} \textit{Alexander Sonderegger}\marginnote{max}

\noindent{Fantasia Contrappuntistica. Edizione definitiva, for piano (1910), Op. 256} \textit{Steven Vanhauwaert}\marginnote{max}

\noindent{Sonatina (No. 1), for piano (1910), Op. 257} \textit{Victor Nicoara}\marginnote{high}

\noindent{Sonatina seconda, for piano (1912), Op. 259} \textit{Victor Nicoara}\marginnote{high}

\noindent{Indianisches Tagebuch. Erstes Buch. [Indian Diary. First Book], for piano (1915), BV 267} \textit{Steven Vanhauwaert}\marginnote{max}

\noindent{Sonatina (no. 3) "ad usum infantis," for piano (1915), BV 268} \textit{Victor Nicoara}\marginnote{high}

\noindent{Sonatina (no. 4) "in diem nativitatis Christi MCMXVII" for piano (1917), BV 274} \textit{Victor Nicoara}\marginnote{high}

\noindent{Sonatina quasi Sonata for piano (1917), BV 275} \textit{Victor Nicoara}\marginnote{high}

\noindent{Sonatina brevis (no. 5), "In Signo Joannis Sebastiani Magni," for piano (1918), BV 280} \textit{Victor Nicoara}\marginnote{high}

\noindent{Kammer-Fantasie über Carmen [Chamber Fantasy after Carmen] (Sonatina no. 6), for piano (1920), BV 284} \textit{Victor Nicoara}\marginnote{high}

\noindent{Toccata: Preludio, Fantasia, and Ciaccona, for piano (1920), BV 287} \textit{Peter Donohoe}\marginnote{max}

\noindent{Drei Albumblätter [Three Albumleaves], for piano (1921), BV 289} \textit{Victor Nicoara}\marginnote{high}
